\section{Finite descent in every rank}\label{sec:finite-descent}

We prove the finite semilocal descent theorem by reducing matrix cocycles to additive ones through successive nilpotent quotients. The key normalization makes every correcting matrix congruent to the identity modulo the radical. It works for arbitrary direct-sum rank and requires no division by an integer.

Write $S=A'$ and $C^q=S^{\otimes_A(q+1)}$. For $h\in S$ put $d(h)=h_2-h_1$ in $C^1$. An additive cocycle $z$ satisfies $z_{23}-z_{13}+z_{12}=0$ in $C^2$; a matrix cocycle $g$ satisfies $g_{12}g_{23}=g_{13}$. A matrix coboundary has the form $h_1^{-1}h_2$. These are the additive and matrix Amitsur cohomology groups and pointed sets of the statement. None of these definitions assumes flatness of $S$ over $A$.

First observe that replacing $A$ by the equalizer $B=H^0(S/A)$ changes none of the tensor powers. Indeed, for $b\in B$ the relation $b\otimes1=1\otimes b$ already holds in $S\otimes_AS$. Thus
\begin{equation}\label{eq:equalizer-tensor}
 S\otimes_AS=S\otimes_BS,
\end{equation}
and the same argument in neighboring slots gives the assertion for all higher tensor powers. In particular the descent categories are identical. The ring $B$ is a finite $A$-module, since it is a submodule of $S$; it is noetherian and semilocal. The finite injective homomorphism $B\to S$ induces a surjection on spectra.

\subsection{The nilpotent calculation}

We establish the equivalence of the two vanishing conditions on every finite initial sequence of radical quotients. The Artinian local factors can be treated separately; a factor on which $S$ is zero contributes nothing. We may therefore take $A_n$ local, with maximal ideal denoted again by $\mathfrak m$, and with $S_n/\mathfrak mS_n$ nonzero. At level zero this is a nonzero finite algebra over the residue field $k$. It is faithfully flat, so ordinary faithfully flat descent gives additive vanishing, triviality of all free-module cocycles, and
\begin{equation}\label{eq:field-equalizer}
 H^0(S_0/k)=k.
\end{equation}

Suppose first that matrix cohomology vanishes in one fixed finite positive rank $r$ at every level. We prove additive vanishing by induction. Given an additive cocycle $z$ at level $n$, subtract the differential of a lift of its primitive at level $n-1$. Its entries now lie in $\mathfrak m^n C_n^1$. This ideal is square-zero, so
\[
 g=(1+z)I_r
\]
is a matrix cocycle. Write $g=h_1^{-1}h_2$. Modulo $\mathfrak m$, the matrix $h$ is invariant, hence is an invertible matrix over $k$ by \eqref{eq:field-equalizer}. Lift it to $\GL_r(A_n)$ and multiply $h$ on the left by its inverse. This preserves the coboundary and arranges $h=I_r+Y$ with entries of $Y$ in $\mathfrak mS_n$. Consequently
\[
 h_2-h_1=z h_1=zI_r,
\]
because $\mathfrak m^{n+1}=0$. The $(1,1)$ entry gives $z=d(h_{11})$. This proves the converse implication even when the fixed rank is one.

Now assume additive vanishing at every level. We prove matrix vanishing in every cardinal rank. For a cardinal $\Lambda$, a matrix means a column-finite matrix acting on the direct sum $S_n^{(\Lambda)}$; its invertibility means that it is an automorphism of that module. Reducing a cocycle modulo $\mathfrak m$ and applying faithfully flat descent over $k$ gives a vector space over $k$. Its dimension is $\Lambda$, as can be checked after tensoring with a residue field of $S_0$. Thus its reduced cocycle is a coboundary.

An invertible column-finite matrix modulo a nilpotent ideal lifts to an invertible matrix. Lift it and its inverse with the same finite supports in each column. The errors in their two products lie in a nilpotent ideal of the endomorphism ring, and finite geometric sums correct them. This ideal is nilpotent uniformly in $\Lambda$: a product of $h$ matrices with entries in an ideal $I$ has entries in $I^h$, since each column involves only finite sums. We can therefore normalize our cocycle as $g=I+X$ with entries in $\mathfrak m C_n^1$.

Suppose inductively that the entries of $X$ lie in $\mathfrak m^q C_n^1$, where $q\geq1$. Modulo $\mathfrak m^{q+1}$ the cocycle equation is additive. The assumed additive vanishing at this quotient supplies entrywise primitives $Y$ with
\[
 dY=X\pmod{\mathfrak m^{q+1}}.
\]
Choose a zero primitive for each zero entry of $X$. The resulting matrix is column-finite. Modulo $\mathfrak m$, each entry of $Y$ is invariant and hence lies in $k$ by \eqref{eq:field-equalizer}. Subtract lifts of these scalars from $A_n$. Their differentials vanish, so the displayed equation is preserved, and now all entries of $Y$ lie in $\mathfrak mS_n$. Thus $h=I+Y$ is invertible. The changed cocycle is $h_1gh_2^{-1}$, and
\[
 h_1g-h_2=X+Y_1-Y_2+Y_1X
             \equiv0\pmod{\mathfrak m^{q+1}}.
\]
Here $Y_1X$ has the required order because $Y$ lies in $\mathfrak m$ and $X$ lies in $\mathfrak m^q$. There is no need to choose $Y$ in $\mathfrak m^q$. Nilpotence finishes this finite iteration. It proves triviality of all cardinal-rank free-module cocycles on every Artinian quotient.

\subsection{Effectiveness over Artinian equalizers}

We explain why the matrix statement proves effectiveness, including the assertion about arbitrary flat modules. A flat module over an Artinian local ring is free in arbitrary rank. To see this, lift a basis modulo the maximal ideal to obtain a map from a free module. Nilpotent Nakayama, valid without finite generation, makes the map surjective. Flatness makes its kernel have zero reduction modulo the maximal ideal; nilpotent Nakayama then makes the kernel zero.

Let $T$ be an Artinian quotient of $A$ and put $S_T=S\otimes_AT$ and $B_T=H^0(S_T/T)$. These are Artinian rings. A flat $S_T$-module is free on each of its finitely many local factors. A descent isomorphism forces its cardinal ranks to agree on factors above the same factor of $B_T$: the tensor product of their residue fields over the residue field of that $B_T$-factor is nonzero, and restriction to any field quotient identifies the two vector-space dimensions. Hence the module has one free cardinal rank on the part above each connected component of $\Spec B_T$.

By \eqref{eq:equalizer-tensor} its datum is a matrix cocycle of that rank. The preceding calculation trivializes it, so the module descends to the free module of the same rank on the corresponding component of $B_T$. A cocycle on one component can be extended by the identity to the others, allowing the proved matrix assertion to be applied component by component. A compatible homomorphism between two such descended modules is a column-finite matrix whose entries belong to the equalizer $B_T$. This proves full faithfulness as well as essential surjectivity.

Conversely, if all flat modules descend effectively, a datum on $S_T^r$ descends to a flat $B_T$-module. On each Artinian local component it is free, and its rank is $r$, as is detected after tensoring with a residue field of $S_T$. Choosing bases makes the original cocycle a coboundary. We have proved all three equivalences in Theorem~\ref{thm:finite-descent}, including their independence of the fixed finite rank.

\subsection{The semilocal ring before completion}

We next deduce the finite-rank assertions over $A$ itself. For an additive cocycle $z\in C^1$, vanishing modulo every power of $\mathfrak m$ says
\[
 z\in\operatorname{im}(d:C^0\to C^1)+\mathfrak m^N C^1
 \quad\text{for every }N.
\]
The quotient $C^1/\operatorname{im}d$ is a finite $A$-module. Krull intersection, since $\mathfrak m\subset\operatorname{rad}A$, makes the intersection of its $\mathfrak m$-powers zero. Thus $z$ is a boundary.

For a finite-rank matrix cocycle $g$, consider the $A$-linear map between finite $A$-modules
\[
 T:M_r(S)\longrightarrow M_r(S\otimes_AS),
 \qquad T(h)=h_2-h_1g.
\]
Artin--Rees provides an integer $c$ such that
\[
 \operatorname{im}T\cap\mathfrak m^{c+1}M_r(S\otimes_AS)
                    \subset\mathfrak m\operatorname{im}T.
\]
Trivialize $g$ modulo a sufficiently large power of $\mathfrak m$ and lift the resulting matrix to $h\in M_r(S)$. Its error $T(h)$ lies in the intersection on the left. There is $w\in\mathfrak mM_r(S)$ with $T(w)=T(h)$, so $h-w$ solves the equation exactly. Its reduction modulo $\mathfrak mS$ is invertible. Since $S$ is finite over the semilocal ring $A$, the ideal $\mathfrak mS$ is contained in the radical of $S$, and $h-w$ is invertible. Hence $g$ is a coboundary. This argument does not require compatible choices of bases on all quotients, or completeness of $A$.

Finally, a finite flat $S$-module has a locally constant finite rank function. Descent identifies the ranks at any two points above the same point of $\Spec B$, by tensoring residue fields as before. This invariant rank function descends to a locally constant function on $\Spec B$: the finite surjective map is closed, so an invariant clopen set and its complement have closed images. Decompose $B$ by the finitely many resulting clopen rank sets. On each piece the module has constant finite rank and is free, since $S$ is semilocal. Matrix vanishing descends its datum to a free $B$-module. The equalizer on matrix entries proves full faithfulness. Thus $\Spec S\to\Spec B$ is effective descent for finite flat modules.

\subsection{Flat base change and arbitrary rank}

Assume now that $A$ is Artinian. Let $R$ be a flat $B$-algebra and put $S_R=S\otimes_BR$. Suppose $M$ is a flat $S_R$-module carrying a descent datum for $S_R/R$. It is flat over $S$, since $S_R$ is flat over $S$. After forgetting the $R$-action, the datum is one for $S/B$: indeed
\[
 (S\otimes_BS\otimes_BR)\otimes_{S\otimes_BR}M
       =(S\otimes_BS)\otimes_SM,
\]
with the appropriate first or second structure map, and the analogous identity holds for the triple tensor product. Artinian effectiveness gives a flat $B$-module $N$ and an isomorphism $S\otimes_BN\simeq M$ compatible with descent.

Multiplication by each $r\in R$ is a morphism of descent data. Full faithfulness descends it uniquely to an endomorphism of $N$. Uniqueness preserves products, sums and the identity, so these endomorphisms make $N$ an $R$-module. The displayed isomorphism is $R$-linear and therefore identifies $S_R\otimes_RN$ with $M$.

It remains to check flatness over $R$. Work on one local factor of $B$, write $\mathfrak b$ for its radical and $k$ for its residue field, and set $R_0=R/\mathfrak bR$ and $S_0=S/\mathfrak bS$. The nonzero finite $k$-algebra $S_0$ is faithfully flat over $k$. Since $M$ is flat over $S_R$, its quotient
\[
 M/\mathfrak bM=(S_0\otimes_kR_0)\otimes_{R_0}(N/\mathfrak bN)
\]
is flat over $S_0\otimes_kR_0$. Faithfully flat descent of flatness makes $N/\mathfrak bN$ flat over $R_0$.

Both $R$ and $N$ are flat over $B$. A free $B$-resolution of $k$, tensored with $R$, is consequently a free $R$-resolution of $R_0$, and gives
\[
 \operatorname{Tor}_i^R(R_0,N)=\operatorname{Tor}_i^B(k,N)=0
 \qquad(i>0).
\]
These facts imply $R$-flatness without any finiteness assumption on $N$ or $R$. For completeness, change of rings gives $\operatorname{Tor}_1^R(Q,N)=0$ for every $R_0$-module $Q$. Every $R$-module has a finite filtration by powers of the nilpotent ideal $\mathfrak bR$, whose quotients are $R_0$-modules. The Tor exact sequences then give the same vanishing for every $R$-module $Q$, which is the flatness criterion.

Full faithfulness over $R$ follows from that over $B$ and the unique descended $R$-actions. This proves effectiveness after every flat affine base change. Affine localization and uniqueness glue the conclusion for an arbitrary flat $B$-scheme. The proof of Theorem~\ref{thm:finite-descent} is complete.
